\documentclass[10pt]{article}

\usepackage[margin=0.8in]{geometry}
\usepackage{mathptmx}
\usepackage{amsmath}
\usepackage{booktabs}
\usepackage{array}
\usepackage[hidelinks,colorlinks=true,linkcolor=blue,citecolor=blue,urlcolor=blue]{hyperref}
\usepackage{enumitem}
\usepackage{titlesec}
\usepackage{parskip}
\usepackage[numbers,sort&compress]{natbib}
\usepackage{xurl}

\setlist{nosep, topsep=2pt, itemsep=2pt}
\titlespacing*{\section}{0pt}{10pt}{4pt}
\titlespacing*{\subsection}{0pt}{8pt}{3pt}

\title{\large Does Voice Compensate When the Face Is Occluded?\\
Audio-Visual Emotion Recognition Under Synthetic Facial Occlusion}
\author{Yashashvi Jangid\\ \small Independent Researcher}
\date{November 2026}

\begin{document}
\maketitle

\begin{abstract}
\noindent
Facial occlusion (masks, sunglasses, hand-over-mouth, random patches) degrades face-based emotion
recognition, but it need not degrade speech. This work asks whether a voice-based classifier can
compensate for occlusion-induced face-branch failure under late fusion, using CREMA-D (7,438 clips,
91 actors, speaker-independent 64/15/12 train/val/test split). Three branches --- eGeMAPS features with
logistic regression, WavLM-base-plus, and a ResNet-18 face branch trained clean or with synthetic
occlusion augmentation --- are combined with four fusion rules (fixed averaging, two confidence-weighted
variants, and a learned gate). Three comparisons were pre-specified before the frozen 12-actor test set
was touched: (P1) fused prediction vs. voice alone, (P2) occlusion-robust vs. clean-trained face
generalization to a held-out occlusion type, and (P3) confidence-weighted vs. fixed fusion. On the test
set, fixed 50/50 fusion of the occlusion-augmented face branch and WavLM reaches 0.778 (95\% CI
[0.707, 0.840]) on clean faces versus 0.705 for WavLM alone, and 0.742 pooled across 12 occlusion
conditions; this fused prediction significantly outperforms voice alone in every one of 13 conditions
(P1). Occlusion-augmented training improves performance on a held-out occlusion type at every severity
(P2). Confidence-weighted fusion shows no improvement over fixed averaging --- pooled, the difference is
statistically detectable but practically negligible ($-0.004$) and favors fixed averaging --- and this is
not uniform across the 13 individual conditions (P3). The result suggests that, at this scale, naive
averaging already captures most of the benefit that explicit reliability weighting was meant to provide.
\end{abstract}

\section{Introduction}
\label{sec:intro}

% AUTHOR TO REWRITE
The motivating question for this work is simple: when a speaker's face is occluded --- by a mask, sunglasses,
a hand, or an arbitrary patch --- does their voice carry enough emotional information to compensate, and
does a fusion rule that explicitly weights branches by their own confidence do a better job of
exploiting that compensation than a fusion rule that does not?

This is a practical question for any audio-visual emotion recognition system deployed outside a
controlled studio: video calls, telehealth, and in-cabin monitoring all involve faces that are sometimes
occluded by ordinary objects, by camera framing, or by the speaker's own hand. A system that degrades
gracefully under occlusion --- by leaning on audio when the face fails --- is more useful than one that does
not.

Three comparisons were specified in writing before the frozen test set was accessed (the pre-specified
plan and its one later amendment are both committed to the project's version-control history; the
amendment's commit hash, \texttt{09e454e981079e7197b46c9c92b0fc2b9ace1a74}, fixes the exact methodology below
before any test number was seen):

\begin{itemize}
\item \textbf{P1.} Does fixed 50/50 fusion of an occlusion-robust face branch and a strong voice branch beat the
  voice branch alone, across clean and occluded conditions?
\item \textbf{P2.} Does a face model trained with synthetic occlusions of three types generalize to a fourth,
  held-out occlusion type (a mask) that it never saw during training?
\item \textbf{P3.} Does confidence-weighted fusion (weighting each branch by its own predicted confidence) beat
  fixed 50/50 averaging, once both branches are temperature-calibrated?
\end{itemize}

P3 is the paper's central methodological question. The intuitive expectation --- and the reason
confidence-weighted fusion is a standard technique in multimodal literature --- is that a branch should be
trusted less when it is uncertain, and an occluded face branch should become visibly less confident. This
work tests that expectation directly and reports whatever the test set shows, including a negative
result.

\section{Related Work}

CREMA-D \citep{cao2014cremad} is a crowd-sourced, single-sentence emotional speech and face dataset with
human recognition accuracies of 40.9\% (audio-only), 58.2\% (visual-only), and 63.6\% (audio-visual),
establishing that fusing the two modalities helps human listeners/viewers on this exact corpus. Several
works fuse audio and visual streams on CREMA-D with temporal or architectural refinements; Birhala,
Ristea, Radoi, and Dutu (2020) \citep{birhala2020temporal} fuse audio and visual features sampled from a temporal window with
independent per-modality offsets and report results on CREMA-D exceeding several prior fusion baselines.

Degraded- and missing-modality robustness has been studied on CREMA-D directly. Goncalves and Busso
(2022a; AuxFormer, ICASSP 2022) \citep{goncalves2022auxformer} combine a transformer with auxiliary single-modality networks and
shared losses, and show the model sustains performance under audio-only and video-only (simulated
missing-modality) conditions on CREMA-D. The journal extension, Goncalves and Busso (2022b; IEEE
Transactions on Affective Computing) \citep{goncalves2022robust}, generalizes this to handle misaligned modalities, weak temporal
structure, and missing features together. Chumachenko, Iosifidis, and Gabbouj (2022; ICPR) \citep{chumachenko2022selfattention} propose a
self-attention fusion architecture evaluated with incomplete audio-visual data and report that dropping
modalities during training improves robustness when a modality is absent or noisy at test time, without
hurting standard-condition performance. Hassanie, Goyal, Ali, and Khan (2026; IEEE Access, vol. 14, pp.
96162--96175) \citep{hassanie2026reliability} study reliability-aware multimodal emotion recognition under heterogeneous and corrupted
audio-visual conditions and report mixed results under video corruption, consistent with the negative
P3 result reported here (Section~\ref{sec:results}).
None of this missing/degraded-modality work that we were able to verify evaluates \emph{deliberately,
systematically occluded faces} across multiple occlusion types and severities with a fixed-vs-confidence-
weighted fusion comparison, which is the specific gap this work addresses; we narrow our claimed
contribution accordingly to: (i) a fixed taxonomy of landmark-placed synthetic occlusion types and
severities, (ii) a pre-specified, statistically tested comparison of confidence-weighted against fixed
fusion under those occlusions, and (iii) evaluation on a frozen, speaker-independent test set reported
exactly once.

Facial occlusion degrades facial expression recognition; Zhang, Verma, Tjondronegoro, and Chandran
(2018) \citep{zhang2018facial} survey this literature from the vision side (occlusion-robust visual features, generative
inpainting, and region-based voting), without addressing audio-visual fusion under occlusion.

\section{Data and Split}

\textbf{Dataset.} CREMA-D contains 7,442 raw audio-video clips from 91 actors; 4 clips are excluded for
documented data-quality reasons (a mislabeled compilation clip and three clips with near-zero or
corrupted video), leaving a canonical set of \textbf{7,438 clips} used throughout. Each clip is a single
spoken sentence labeled with one of six emotions (anger, disgust, fear, happiness, neutral, sadness).

\textbf{Split.} Actors are partitioned into training (64 actors), validation (15 actors), and test (12
actors) sets, speaker-independent --- no actor appears in more than one split. The split was generated with a
fixed random seed (seed 42, sex$\times$race-stratified with per-stratum largest-remainder rounding) and
was frozen before any model development began; it is enforced by an automated leakage check on every
data load. The 12-actor test set was not touched by any model training, hyperparameter selection, or
calibration decision; it was accessed exactly once, for scoring only, after the evaluation plan below was
committed, and the committed record of that single run is at commit
\texttt{d2af74c65605b4b81b9a4333ce81c597115a6a7e}.

\section{Methods}

\textbf{Audio branch 1 --- eGeMAPS + logistic regression.} 88-dimensional eGeMAPSv02 functional acoustic
features \citep{eyben2016gemaps} (extracted with openSMILE) are standardized and classified with a multinomial logistic
regression, with its regularization strength tuned on the validation set.

\textbf{Audio branch 2 --- WavLM.} \texttt{microsoft/wavlm-base-plus} \citep{chen2022wavlm}, with its convolutional feature encoder frozen
and a small two-layer classification head fine-tuned on top of a masked mean pool over the transformer's
final hidden states, trained on the 64 training actors with early stopping on validation macro-F1, 3
seeds. Inference uses a batch size of one: the frozen convolutional encoder has no attention-masking
mechanism, so zero-padding a batch of variable-length clips affects the encoder's outputs, meaning
predictions depended on batch composition --- a batch-size-of-one inference pass removes this dependency
so that a clip's score depends only on its own audio and the trained weights.

\textbf{Face branch --- ResNet-18.} A ResNet-18 \citep{he2016resnet} image classifier is fine-tuned on face crops extracted from
video frames via MediaPipe face landmarking, in three variants: \textbf{clean-trained} (no occlusion during
training), \textbf{aug\_all} (all four synthetic occlusion types applied during training), and \textbf{aug\_no\_mask}
(three occlusion types applied during training, with the mask type withheld so that its test-time
performance on masks reflects generalization from the other three occlusion types rather than exposure
to the mask generator itself; Section~\ref{sec:limitations} notes a caveat on how cleanly ``unseen'' this is in practice).
Clip-level prediction averages per-frame softmax probabilities across multiple sampled frames per clip.

\textbf{Synthetic occlusions.} Four occlusion types are rendered on top of a detected face using its
MediaPipe landmarks: a lower-face \textbf{mask} (coverage increases continuously with severity, from
mouth-and-chin at the lowest severity to just below the eyes at the highest), \textbf{sunglasses} (an opaque
ellipse over the eye region), a randomly placed gray \textbf{patch}, and a \textbf{hand-over-mouth} occlusion
anchored at the mouth landmarks. Each type is applied at three severities, giving 12 occlusion
conditions plus the clean condition, 13 in total.

\textbf{Fusion.} Four methods combine the two per-branch probability vectors: \textbf{fixed 50/50} averaging;
\textbf{confidence-weighted (max-probability)}, which weights each branch by its own calibrated maximum
softmax probability; \textbf{confidence-weighted (negative entropy)}, an analogous weighting using negative
entropy as the confidence signal; and a \textbf{learned gate}, a logistic regression predicting whether the
face branch is correct from both branches' confidence scores, used as the face branch's fusion weight.
Both per-branch probability vectors are temperature-calibrated before fusion. Per the pre-specified
amendment, both the calibration temperatures and the learned gate are fit once on the 15 validation
actors and frozen; the 12 test actors' labels are used only to score predictions, never to fit any
calibration or gating parameter.

\section{Evaluation Protocol}
\label{sec:evalprotocol}

\textbf{Pre-specified evaluation plan.} A written evaluation plan specifying the exact models, conditions,
fusion methods, and a set of statistical comparisons was committed to version control before the test
set was accessed. In that original plan, the two primary comparisons used the eGeMAPS audio branch (the
only audio branch trained at the time); the fusion vs. audio-only comparison was restated in terms of
WavLM, and P1--P3 as defined in Section~\ref{sec:intro} were adopted as primary, in a single amendment made and
committed (hash \texttt{09e454e981079e7197b46c9c92b0fc2b9ace1a74}) once WavLM was available and before the test
actors were accessed. That amendment also froze calibration and the fusion gate on validation data
rather than re-fitting them on the test actors. No further changes were made after test results were
observed. The original eGeMAPS-based comparisons are still reported, as secondary results (Table~\ref{tab:secondary}).

\textbf{Metric.} Macro-averaged F1 (macro-F1) across the six emotion classes is the primary metric throughout,
reported alongside accuracy. Every model involving stochastic training (the face branch, WavLM) is
trained with 3 random seeds and reported as a mean $\pm$ standard deviation across seeds; per-seed scores for
every number in this report are available in the project's committed results log (see Section~\ref{sec:availability}).
Statistical comparisons use an actor-level paired bootstrap: the resampling unit is the actor (not
the clip), because clips from the same actor are not independent observations. Each of 2,000 bootstrap
draws resamples the 12 test actors with replacement, pools all of a drawn actor's clips, computes the
paired macro-F1 difference between the two methods being compared for each of the 3 seeds using that same
resampled actor set, and averages the difference across seeds. The reported 95\% confidence interval is
the 2.5th/97.5th percentile of this distribution of 2,000 seed-averaged differences; the two-sided
p-value is $p = 2 \times \min(\Pr[\text{diff} \le 0],\ \Pr[\text{diff} \ge 0])$, i.e. twice the smaller of
the two one-sided bootstrap tail proportions. We use \textbf{``pre-specified''} rather than ``pre-registered''
throughout this report, since the commitment mechanism was version-control history on a private project
repository rather than a registry with external timestamping.

\textbf{Abstention rule.} The face branch can fail to detect a face at all on a heavily occluded frame. A
clip where every sampled frame fails detection is scored as an abstention: it counts as a false negative
for its true emotion class, is never counted as a (possibly wrong) prediction of any class, and remains
in the accuracy denominator. Neither audio branch (eGeMAPS+LR nor WavLM) abstains.

\section{Results}
\label{sec:results}

All results below are from the frozen 12-actor test set, evaluated exactly once. Pooled figures average
over the 12 occlusion conditions (clean excluded); per-condition breakdowns are given for the three
pre-specified comparisons.

\begin{table}[h]
\small
\centering
\caption{Branch-level macro-F1 (test set, mean $\pm$ std over 3 seeds, with actor-level bootstrap 95\% CI
where computed).}
\label{tab:branch}
\begin{tabular}{llcc}
\toprule
Branch & Condition & Macro-F1 & 95\% CI \\
\midrule
eGeMAPS + LR (tuned) & --- (deterministic) & 0.5044 $\pm$ 0.0000 & [0.4576, 0.5424] \\
WavLM & --- & 0.7049 $\pm$ 0.0121 & [0.6269, 0.7616] \\
Face, clean-trained & clean & 0.6050 $\pm$ 0.0151 & [0.5034, 0.6987] \\
Face, aug\_all & clean & 0.6227 $\pm$ 0.0100 & [0.5207, 0.7143] \\
Face, aug\_no\_mask & clean & 0.6303 $\pm$ 0.0137 & [0.5196, 0.7279] \\
\bottomrule
\end{tabular}
\end{table}

\begin{table}[h]
\small
\centering
\caption{Fusion (WavLM + face), macro-F1 mean over 3 seeds, pooled over the 12 occlusion conditions and
on the clean condition alone.}
\label{tab:fusion}
\begin{tabular}{lcc}
\toprule
Face variant & Pooled (12 occl.) & Clean \\
\midrule
clean-trained & 0.710 & 0.7754 $\pm$ 0.0163 \\
aug\_all & 0.742 & 0.7781 $\pm$ 0.0095 \\
aug\_no\_mask & 0.739 & 0.7781 $\pm$ 0.0120 \\
\bottomrule
\end{tabular}
\end{table}

Fixed 50/50 fusion is the best of the four fusion methods for every face/audio-branch pairing tested
except clean-trained face + eGeMAPS, where the learned gate edges it out by 0.003 macro-F1.

\textbf{P1 --- fixed fusion vs. WavLM alone (aug\_all face).} Significant in every one of the 13 conditions.
Pooled: $+0.0404$ (95\% CI [0.0182, 0.0629], $p = 0.001$). The effect is largest on clean faces ($+0.0731$,
CI [0.0355, 0.1072]) and shrinks under the heaviest occlusion (mask 0.9: $+0.0099$; patch 0.9: $+0.0108$) but
never loses significance or reverses sign.

\textbf{P2 --- aug\_no\_mask vs. clean-trained face-only, under the held-out mask type.} Significant at every
severity: $+0.0935$ (severity 0.3, CI [0.0496, 0.1331]), $+0.2258$ (severity 0.6, CI [0.1802, 0.2742]), and
$+0.0584$ (severity 0.9, CI [0.0400, 0.0766]), all $p < 0.001$.

\textbf{P3 --- confidence-weighted (max-prob) vs. fixed fusion (aug\_all + WavLM).} Pooled across the 13
conditions --- the pre-specified test --- confidence-weighted fusion shows \textbf{no improvement over fixed
averaging}: the difference is $-0.0038$ (95\% CI [$-0.0050$, $-0.0025$], $p < 0.001$), statistically detectable
but practically negligible in size, and in fixed averaging's favor. The 13 individual conditions were
not themselves pre-specified as separate tests and are reported here as a \textbf{descriptive} breakdown of
the pooled result, not as 13 independent confirmatory tests: 8 reach the conventional 0.05 threshold in
fixed averaging's favor --- mask at severities 0.3 ($p=0.012$), 0.6 ($p=0.045$), 0.9 ($p=0.008$); sunglasses
at 0.3 ($p<0.001$), 0.6 ($p=0.021$), 0.9 ($p<0.001$); and patch at 0.3 ($p<0.001$) and 0.6 ($p=0.001$) ---
while 4 do not: patch 0.9 ($p=0.830$) and hand-over-mouth at 0.3 ($p=0.093$), 0.6 ($p=0.557$), and 0.9
($p=0.069$); clean lands exactly at the threshold ($p=0.050$). That is $8+4+1=13$ conditions in
total.

\begin{table}[h]
\small
\centering
\caption{Secondary comparisons: fixed fusion vs. audio-only, by audio branch and face variant, pooled
over all 13 conditions (seed-averaged paired bootstrap).}
\label{tab:secondary}
\begin{tabular}{llccc}
\toprule
Audio branch & Face variant & Diff & 95\% CI & $p$ \\
\midrule
eGeMAPS & clean-trained & $+0.0311$ & [$-0.0125$, 0.0753] & 0.169 \\
WavLM & clean-trained & $+0.0098$ & [$-0.0098$, 0.0263] & 0.336 \\
eGeMAPS & aug\_all & $+0.1221$ & [0.0767, 0.1681] & $<0.001$ \\
WavLM & aug\_all (= P1) & $+0.0404$ & [0.0182, 0.0629] & 0.001 \\
eGeMAPS & aug\_no\_mask & $+0.1085$ & [0.0582, 0.1595] & 0.001 \\
WavLM & aug\_no\_mask & $+0.0366$ & [0.0122, 0.0596] & 0.002 \\
\bottomrule
\end{tabular}
\end{table}

These comparisons were the project's original primary tests before the amendment described in Section~\ref{sec:evalprotocol}
moved the primary audio branch to WavLM; they are reported here as secondary results. Fusion only
reliably beats audio-only, for either audio branch, once the face branch is occlusion-robust
(aug\_all/aug\_no\_mask); with the clean-trained face, neither audio branch's fusion-vs-audio-only
comparison reaches significance pooled.

\section{Discussion}

\textbf{Why fixed averaging is hard to beat.} When one branch's output distribution is close to uniform
(the behavior of an uncertain or abstaining classifier), averaging it with a confident branch barely
perturbs the confident branch's peak class, because a near-uniform vector adds roughly the same small
amount to every class and so does not change the arg-max. Explicit confidence weighting earns its keep
when a branch is confidently wrong rather than merely hedged --- on development (validation) data, the
clean-trained face branch's output under heavy masking was observed to remain relatively peaked (less
confident than on clean faces, but still more confident than its accuracy under masking would justify)
rather than hedging toward uniform --- but this is the exception rather than the rule across the 13
conditions tested, and the machinery needed to detect it (temperature scaling, a learned gate, or an
entropy-based signal) introduces its own variance. Fixed averaging has no corresponding failure mode,
because it never acts on either branch's self-reported confidence at all.

\textbf{The P3 result is the paper's central finding, not an afterthought.} The motivating hypothesis behind
P3 was intuitive enough to be worth testing formally rather than assuming: a system should trust a branch
less when that branch is less sure of itself, and an occluded face branch should become visibly less
confident. The pre-specified test shows that --- at this scale of training data, with these two specific
audio and face architectures, and with the particular confidence signals tested (calibrated max-softmax
probability, negative entropy, and a learned logistic gate) --- this intuition does not translate into a
practical improvement: confidence weighting shows no improvement over fixed averaging pooled, and fixed
averaging is a strictly simpler, equally or more effective default here. One possible interpretation of
where the (descriptive, not independently tested) per-condition gap is largest and smallest: under
sunglasses and mid-severity patch, the face branch may be degraded but not yet collapsed to a uniform or
systematically wrong output, whereas under hand-over-mouth and severe patch it may have collapsed into a
regime --- uniform abstention, where fixed and confidence-weighted fusion converge by construction, or some
other regime --- that this project's confidence signals do not distinguish from genuine uncertainty. This
reading is plausible given the per-condition pattern but was not itself tested and should be treated as
an interpretation, not a finding.

\section{Limitations}
\label{sec:limitations}

\begin{itemize}
\item \textbf{Synthetic occlusions.} All occlusions are rendered digitally on top of clean footage using detected
  facial landmarks, not captured from genuinely occluded recordings. Real occlusions (an actual surgical
  mask, real sunglasses, a real hand) vary in material, fit, lighting interaction, and motion in ways a
  landmark-anchored synthetic overlay cannot fully reproduce.
\item \textbf{Occlusions are placed using landmarks detected on the clean, unoccluded face}, not re-detected on an
  already-occluded image. A real deployment would need to detect facial landmarks on a face that may
  already be occluded, which is a harder and potentially less reliable detection problem than the one
  this project's synthetic-occlusion pipeline solves; this project's face-detection failure handling
  (the abstention rule, Section~\ref{sec:evalprotocol}) applies to the downstream emotion-classification step, not to this
  upstream landmark-placement step.
\item \textbf{aug\_no\_mask's held-out occlusion is not fully independent of the occlusion types seen in
  training.} The hand-over-mouth occlusion (seen during aug\_no\_mask training) and the mask occlusion
  (withheld) both cover overlapping lower-face regions. aug\_no\_mask's improvement on masks at test time
  may therefore partly reflect transfer from a visually similar occlusion region rather than transfer
  from occlusion robustness in general; we describe this as training-time exclusion of one named
  occlusion type, not as a test of robustness to an occlusion region the model never saw in any form.
\item \textbf{Acted speech.} CREMA-D emotions are performed by actors reading fixed sentences, not spontaneous
  affect. Results here characterize performance on acted, scripted emotional expression and may not
  transfer directly to spontaneous or naturalistic emotional speech and expression.
\item \textbf{A single dataset.} All results are from CREMA-D; no cross-corpus evaluation was performed, so it is
  unknown how these findings --- particularly the P3 null result for confidence weighting --- generalize to
  other audio-visual emotion corpora, recording conditions, or languages.
\item \textbf{A 12-actor test set and correspondingly wide confidence intervals.} Every reported test-set
  confidence interval reflects the actor-level sampling uncertainty of a 12-actor set; several absolute
  (non-paired) intervals overlap each other even where the corresponding \emph{paired} comparison was
  significant (e.g., fixed-fusion and WavLM-alone absolute intervals on clean faces overlap, even though
  P1's paired test on the same data is significant). Paired comparisons, which exploit the correlation
  between two methods' scores on the same resampled actors, are more sensitive than comparing two
  independent intervals and are the basis for every significance claim in this report.
\item \textbf{WavLM training-time batch-padding noise.} The batch-composition sensitivity addressed at inference
  time (by scoring one clip per forward pass) was also present during WavLM's own fine-tuning on Kaggle,
  where batches were padded and shuffled every epoch --- an uncharacterized source of training stochasticity
  beyond the usual data-order and weight-init randomness. This does not invalidate the checkpoints used
  here; retraining with length-bucketed batching to reduce it was not attempted in this project.
\item \textbf{A demographic age difference observed only on development (validation) data, not tested on the
  frozen test set, and not attributed to any cause.} A secondary, single-training-seed analysis on the
  79 non-test actors, using the eGeMAPS audio branch (not WavLM), found an \emph{observed} macro-F1 difference
  --- its 95\% bootstrap confidence interval excludes zero, which we call confirmed only in the narrow
  statistical sense that the interval excludes zero at this sample size, not confirmed as a real-world
  effect with a known cause --- between actors aged 30 or under and actors aged 46--60, for the
  occlusion-augmented face branch ($+0.134$, 95\% CI [0.068, 0.201]) and the corresponding fusion pipeline
  ($+0.095$, CI [0.022, 0.170]); the audio branch alone showed no such difference ($+0.004$, CI
  [$-0.057$, 0.068]). A further, paired analysis asked whether occlusion-augmented training widens this age
  difference relative to clean training: for the fusion pipeline, the widening itself was also an observed
  difference with a CI excluding zero ($+0.039$, 95\% CI [0.006, 0.075]); for the face branch alone the
  corresponding widening did not reach that threshold ($+0.057$, CI [$-0.020$, 0.125]). This analysis used 1
  training seed per fold (a reported scope reduction from the usual 3) and was deliberately kept out of
  scope for the frozen test-set evaluation; no cause for either difference was investigated or is claimed.
  It is reported as a preliminary, development-stage observation for a future, adequately powered study to
  test formally, not as a claim about the test set or about why any difference exists.
\item \textbf{Test-set scores are higher than validation-set scores for every model and condition compared.} Both
  audio branches, all three face variants, and the fusion pipeline score somewhat higher on the 12-actor
  test set than the 15-actor validation set (e.g., WavLM: 0.6929 validation vs. 0.7049 test; clean-trained
  face on clean faces: 0.5724 validation vs. 0.6050 test). Reported as observed; with only 12 and 15
  actors respectively this is consistent with ordinary actor-sampling variability between two disjoint
  speaker sets, not evidence of a systematic split difference, and was not investigated further as it was
  outside this evaluation's pre-specified scope.
\end{itemize}

\section{Conclusion}

% AUTHOR TO REWRITE
On a frozen, speaker-independent test set of 12 actors, fusing a strong voice classifier with an
occlusion-robust face classifier significantly outperforms voice alone across every occlusion condition
tested, including the clean condition (P1); training the face branch on synthetic occlusions that
exclude a given occlusion type still measurably improves performance on that excluded type at test time
(P2, with the caveat on occlusion-region overlap noted in Section~\ref{sec:limitations}); and, contrary to the intuitive
expectation that motivated much of this project's earlier development, explicitly weighting each branch
by its own confidence shows no improvement over fixed averaging pooled ($-0.004$, statistically detectable
but practically negligible, and in fixed averaging's favor), and this is not uniform condition-by-condition
(P3). The practical implication for a deployed audio-visual emotion
recognition system handling occluded faces is that a simple, parameter-free fusion rule is a strong and
difficult-to-beat default; the engineering effort involved in building and calibrating a
confidence-weighted or learned fusion mechanism was not justified by this project's results, on this
dataset, at this scale.

\section{Data and Code Availability}
\label{sec:availability}

Code, configuration, the pre-specified evaluation plan and its amendment, and the full results log
(\texttt{RESULTS.md}) referenced throughout this report are available at
\url{https://github.com/yashashvi422005/multimodal-emotion-recognition}. CREMA-D \citep{cao2014cremad} is distributed under the Open
Database License (ODbL) v1.0, with its contents under the Database Contents License (DbCL) v1.0; it is
not redistributed by this project.

\section*{Acknowledgements (Use of AI Assistants)}

This project's code --- data processing, model training, evaluation, calibration, bootstrap statistics,
and the Kaggle notebook used for WavLM fine-tuning --- was developed with substantial assistance from
Claude Code (Anthropic), an AI coding assistant, operating under the direction and review of the author
at every step. Claude in a separate conversational interface was used for planning and for review of
methodological decisions (the split protocol, the calibration-freezing amendment, the choice of
pre-specified comparisons, and scope reductions) alongside the author. Every such decision was made or
explicitly approved by the author before being implemented or written. Claude Code was also used to draft
this report's prose from the project's committed, author-reviewed results; all numbers in this report
were read directly from the project's version-controlled results log rather than computed by, or left to
the discretion of, the assistant while drafting. The author takes full responsibility for the correctness
and interpretation of all content in this report.

\bibliographystyle{unsrtnat}
\bibliography{references}

\end{document}
